% STATE: DRAFT (awaiting human review).
\section{Limitations and Future Work}\label{sec:limitations}

\subsection{Limitations}

% PROV[user]: Scope discussion: finite LTL checker, explicit seriality and
% no deadlock completion; CTL* deferred. Earlier learning discussion: model guarantees.
% PROV[code]: KILT/Kripke/Check.lean:L14-L44;
% KILT/Kripke/PreparedCheck.lean:L19-L31.
The checker handles LTL properties of finite Kripke structures with explicit
initial states and a separate seriality proof. Deadlocked models remain
representable, but the public checking contract requires every state to have
a successor. Seriality supplies continuation; it does not impose fairness.
The correctness results concern the declared transitions, labels and trace
semantics. Their application to a particular system therefore depends on the
chosen model describing that system appropriately.

% PROV[user]: 16 November challenges discussion: larger stress tests needed;
% VM resources and time restricted evaluation.
% PROV[code]: KILT/LTL/Compile.lean, card_automatonState_le;
% report/src/sections/05-examples-and-evaluation.tex, experimental-cost subsection.
Tableau state enumeration is exponential in the number of subformulas, and
product construction combines automaton states with system states.
Reachable-state preparation and cached successor rows reduce work, but larger
models still need stress testing. The VM resource limits and time constraints
described in \cref{sec:learning} restricted this part of the evaluation;
\cref{sec:experimental-cost} records the measurements and their scope.

% PROV[user]: Current request: explain failure conditions and what to do when
% metaprogramming fails; prior API-friction discussion requested a lightweight API.
% PROV[code]: KILT/Frontend.lean:L14-L46;
% Examples/Demos/Frontend.lean, symbolic_manual and final computational-cost note.
\paragraph{When automation fails.}
The macros target small, concrete goals. A failure of \lean{kilt_verify}
does not itself establish that a property is false: the labels may be
symbolic, computation may be too expensive, or the property may actually
fail. The macros currently expose errors from their
underlying tactics. For symbolic goals, an explicit semantic proof can use
the available hypotheses.
\sourcefile{Examples/Demos/Frontend.lean} demonstrates this with
\lean{symbolic_manual}. For larger concrete models, the execution API can
compute a verdict and extract a lasso; printed output is distinct from a
proof of that concrete verdict. After a negative check, the prepared API
provides the counterexample for inspection. The verdict and witness contracts
are described in \cref{sec:correctness}.

\subsection{Future Work}

% PROV[user]: Current request: seek Zulip review and feedback; develop a serious
% library, possibly on a new branch; files too fragmented for assignment clarity.
\textbf{Review and library structure.}
A next step is to seek feedback on Zulip from the Lean and CSLib communities,
covering the definitions, theorem interfaces and project organisation. A
separate development branch would allow that feedback to guide the transition
from an assignment to a maintained library. The fine module boundaries helped
present individual stages, but some files now feel unnecessarily fragmented.
Consolidation should make related results easier to find while keeping the
main subsystem interfaces clear.

% PROV[user]: Current request and earlier scope discussion: pursue CSLib
% contributions from the independent generalized-Buchi, tableau and graph layers.
% PROV[code]: KILT/Automata/GeneralizedBuchi.lean;
% KILT/Automata/GeneralizedBuchiCorrect.lean; KILT/Automata/RoundRobin.lean;
% KILT/LTL/TableauCorrect.lean; KILT/Graph/Prepare.lean;
% KILT/Automata/ProductLasso.lean.
\textbf{Contributions as separate components.}
KILT contains several components that can be reviewed independently:
generalised B\"uchi automata and their round-robin conversion, the LTL tableau
construction, and finite graph search with lasso extraction. These provide
concrete starting points for proposed CSLib contributions, subject to review
and adaptation to the upstream interfaces and conventions.

% PROV[user]: Current request: metaprogramming failure handling; earlier frontend
% discussion and .ai/human/MAYBE_BETTER.md, section 11, user-centred error reporting.
\textbf{Frontend and evaluation.}
Further metaprogramming work should make failure easier to interpret:
distinguish an unsuccessful proof attempt from a computed negative verdict,
and direct the user towards an explicit proof or counterexample inspection
as appropriate. This would improve the existing thin frontend. Larger stress
tests, with more time and suitable hardware, would also give a better basis
for choosing subsequent caching and search improvements.

% PROV[user]: Current request and clarification: CTL*, not LTL*; improve
% under-documented parts. Earlier scope discussion explicitly deferred CTL*.
\textbf{Broader scope and documentation.}
CTL* remains a longer-term extension beyond the current LTL toolkit.
Improving the explanations of under-documented modules is also part of the
library work, especially the assumptions, contracts and connections between
the independently reusable components.

\clearpage
